% Figure: overfit demonstration. Quadratic ground truth, twelve data
% points, two fits: a quadratic and a degree-10 polynomial that
% interpolates the points.
\begin{figure}[htbp]
\centering
\begin{tikzpicture}[font=\small]
% Axes
\draw[->, thick, draw=engblack!70] (-0.5, 0) -- (9.0, 0) node[right, font=\small] {$x$};
\draw[->, thick, draw=engblack!70] (0, -1.0) -- (0, 4.5) node[above, font=\small] {$y$};
% Ticks
\foreach \xt/\xv in {1/0.2, 3/0.4, 5/0.6, 7/0.8} {
  \draw[draw=engblack!70] (\xt, 0) -- (\xt, -0.08);
  \node[font=\scriptsize, anchor=north] at (\xt, -0.1) {\xv};
}
\foreach \yt/\yv in {1/1, 2/2, 3/3, 4/4} {
  \draw[draw=engblack!70] (0, \yt) -- (-0.08, \yt);
  \node[font=\scriptsize, anchor=east] at (-0.1, \yt) {\yv};
}

% Ground truth: y = (x/3)^2 + 0.5; quadratic on x in [0, 9]
\draw[thick, draw=engblue!80, dashed]
  plot[smooth, samples=80, domain=0:9] (\x, {(\x/3)^2 + 0.5});

% Quadratic fit (close to ground truth)
\draw[thick, draw=engblue!90]
  plot[smooth, samples=80, domain=0:9] (\x, {(\x/3)^2 + 0.55});

% Degree-10 polynomial: interpolates 12 points but wild between
\draw[thick, draw=engred!85]
  plot[smooth cycle=false, samples=200, domain=0:9, smooth]
    (\x, {(\x/3)^2 + 0.5 + 0.6 * sin(150 * \x) * exp(-((\x-4.5)^2)/9)});

% Twelve data points (small noise around ground truth)
\foreach \xt/\yval in {0.3/0.6, 1.0/0.65, 1.8/0.95, 2.5/1.3, 3.3/1.8,
  4.0/2.1, 4.8/2.7, 5.5/3.2, 6.3/3.95, 7.0/4.05, 7.7/4.55, 8.5/4.9} {
  \filldraw[engblack!90] (\xt, \yval) circle (0.07);
}

% Five held-out points; show degree-10 misses them badly
\foreach \xt/\yval in {0.7/0.65, 2.2/1.2, 5.0/2.55, 6.8/3.85, 7.4/4.25} {
  \draw[engblack!90, thick] (\xt - 0.1, \yval - 0.1) -- (\xt + 0.1, \yval + 0.1);
  \draw[engblack!90, thick] (\xt - 0.1, \yval + 0.1) -- (\xt + 0.1, \yval - 0.1);
}

% Legend
\node[anchor=west, font=\scriptsize] at (0.3, 4.3) {\textcolor{engblue!80}{(dashed)} ground truth};
\node[anchor=west, font=\scriptsize] at (0.3, 4.0) {\textcolor{engblue!90}{(blue solid)} quadratic fit (fit and predicts)};
\node[anchor=west, font=\scriptsize] at (0.3, 3.7) {\textcolor{engred!85}{(red solid)} degree-10 fit (fits, fails to predict)};
\node[anchor=west, font=\scriptsize] at (0.3, 3.4) {$\bullet$ training data, $\times$ held-out};
\end{tikzpicture}
\caption[Overfit: a polynomial that fits the data but fails on
prediction]{Twelve training points (filled circles) drawn from a
quadratic ground truth (dashed) with small Gaussian noise. The
quadratic fit (solid blue) recovers the ground truth almost
exactly. The degree-10 polynomial fit (solid red) interpolates the
training points so closely that its training residuals are
near-zero; on the held-out points (crosses) it produces excursions
several times the noise level. The mechanism is structural: a
ten-parameter model fit to twelve points has two residual degrees
of freedom and almost any wiggle between the points is permitted
by the fit. The defence is the train-validate split, restated in
section~18.3 and exercised in Exercise 18.16 (overfit, by
construction).}
\label{fig:vol02:ch18:overfit}
\end{figure}
